%FILE IN LATEX DOVE VERRÀ SCRITTA LA RELAZIONE
\documentclass[10pt]{article}
\usepackage[utf8]{inputenc}
\usepackage{amsmath}
\usepackage{amsfonts}
\usepackage{amssymb}
\usepackage{graphicx}
\usepackage{hyperref}
\usepackage{geometry}
\usepackage{pgfplots}

\pgfplotsset{compat=1.18}
\geometry{a4paper, margin=1in}
\title{Analisi delle Prestazioni di Inserimento in Alberi Binari di Ricerca}
\author{
    Londero Stefano\\
    \small{172585}
    \and
    Grasso Federico\\
    \small{173731}
}

\date{Algoritmi e Strutture Dati - 2026}

\begin{document}
\maketitle

\section*{Introduzione}
L'obiettivo del progetto è confrontare tra Alberi Binari di Ricerca (BST), Alberi (AVL) e Red-Black (RB) il tempo medio di inserimento di una chiave quando l'albero contiene n chiavi e, ottenere per diversi valori di n, stime robuste dei tempi di esecuzione.

\section*{Richiami Teorici}
Il costo dell'inserimento in un albero di ricerca dipende dalla lunghezza del cammino radice-foglia, ed è quindi proporzionale alla sua altezza $h$.
\begin{itemize}
    \item \textbf{BST Semplice:} Non garantisce alcun bilanciamento. Nel caso peggiore può degenerare in $\mathcal{O}(n)$, portando a tempi di inserimento lineari.

    \item \textbf{Alberi AVL:} Mantengono un bilanciamento stretto garantendo che la differenza di altezza tra i sotto-alberi di ogni nodo sia al massimo 1. Nel caso peggiore vengono assicurati tempi asintotici $\mathcal{O}(\log n)$ al costo di frequenti rotazioni strutturali.

    \item \textbf{Alberi Red-Black:} Utilizzando colorazioni per i nodi viene garantita una complessità $\Theta(\log n)$. Sebbene il cammino massimo possa essere leggermente più lungo rispetto a un AVL, il RBT richiede tipicamente meno operazioni di rotazione per ripristinare le proprietà dopo un inserimento.
\end{itemize}

\section*{Metodologia}
Per garantire misurazioni robuste e stabili, il progetto si basa su due pilastri:
\subsubsection*{Varietà Morfologica}
Invece di 
\begin{enumerate}
    \item Selezionare una chiave non ancora presente e se ne misura l'inserimento tramite \texttt{time.perf\_counter()}.
    \item Estrarre casualmente una chiave dall'albero e la si rimuove \textit{senza cronometrare} l'operazione.
\end{enumerate}
Questa tecnica garantisce che prima di ogni inserimento (e dunque prima che il cronometro parta), la dimensione torna sempre a essere esattamente $n$. Inoltre la continua eliminazione e reinserimento di nodi altera in modo casuale la topologia della struttura dati, garantendo che le misurazioni avvengano su albero con forme uniformemente distribuite.
\subsubsection*{Struttura del Progetto}
Il progetto è stato strutturato nel seguente modo, seguendo un'ottica di programmazione ad oggetti:
\begin{itemize}
    \item \textbf{Codici delle Strutture Dati}: I file \texttt{BST.py}, \texttt{AVL.py} e \texttt{RBT.py} contengono i codici completi delle strutture dati. In particolare nei codici delle ultime due strutture dati il codice è stato alleggerito riciclando alcune parti del codice da \texttt{BST.py}.
    \item \textbf{KeyManager}: Questo file implementa un partizionamento di un array, garantendo la selezione di chiavi libere e occupate in un tempo 
\end{itemize}


\section*{Misurazione}
Per la misurazione abbiamo utilizzato metodi come, per esempio, \texttt{perf\_counter()} per misurare il tempo di esecuzione dell'inserimento di chiavi in alberi BST, AVL e RB. Abbiamo eseguito più iterazioni per ogni valore di $n$ per ottenere stime robuste dei tempi di esecuzione. Abbiamo anche utilizzato grafici per visualizzare i risultati e confrontare le prestazioni dei diversi tipi di alberi. I grafici mostrano il tempo medio di inserimento in funzione del numero di chiavi presenti nell'albero, evidenziando le differenze tra BST, AVL e RB.

\subsection*{Algoritmo}
Il progetto implementa un sistema di misurazione delle prestazioni per tre diverse strutture dati ad albero: Alberi Binari di Ricerca (BST), Alberi AVL e Alberi Rosso-Nero (RBT). L'obiettivo è quantificare il tempo di inserimento di una chiave in ciascun tipo di albero al variare della dimensione dell'albero stesso.

\subsection*{Struttura del Progetto}
Il progetto è organizzato nei seguenti moduli:
\begin{itemize}
    \item \textbf{ExperimentSetup}: Configura i parametri dell'esperimento, calcolando una progressione geometrica di valori di $n$ da testare (da 1000 a 10.000.000 elementi) e inizializzando gli alberi con $n$ chiavi casuali.
    \item \textbf{ExperimentRunner}: Esegue una serie di esperimenti controllati per ciascun tipo di albero e per ogni valore di $n$.
    \item \textbf{BST, AVL, RBT}: Implementano le tre strutture dati con operazioni di inserimento, ricerca e rimozione.
    \item \textbf{KeyManager}: Gestisce l'insieme di chiavi utilizzate negli esperimenti.
    \item \textbf{Plotter}: Visualizza i risultati degli esperimenti tramite grafici.
\end{itemize}

\subsection*{Algoritmo di Misurazione}
Per ogni valore di $n$ nella progressione geometrica, l'algoritmo esegue i seguenti passi:
\begin{enumerate}
    \item Inizializzazione: Viene creato un albero vuoto (BST, AVL o RBT).
    \item Popolazione iniziale: L'albero viene popolato con $n$ chiavi casuali generate da una lista predefinita di chiavi.
    \item Ciclo di misurazione: Per 10 iterazioni (configurabili):
    \begin{enumerate}
        \item Viene selezionata una chiave disponibile dall'insieme delle chiavi non ancora inserite.
        \item Si misura il tempo di inserimento di questa chiave nell'albero utilizzando \texttt{time.perf\_counter()}.
        \item Viene rimossa casualmente una chiave dall'albero per mantenere la dimensione costante e simulare un carico dinamico.
    \end{enumerate}
    \item Calcolo della mediana: Viene calcolata la mediana dei 10 tempi di inserimento misurati, che rappresenta una stima robusta delle prestazioni.
\end{enumerate}

\subsection*{Caratteristiche Implementative}
\begin{itemize}
    \item \textbf{Progressione geometrica}: I valori di $n$ sono generati secondo una progressione geometrica per coprire uniformemente lo spazio logaritmico tra il minimo e il massimo.
    \item \textbf{Bilanciamento dinamico}: AVL e RBT mantengono automaticamente il bilanciamento dell'albero dopo ogni inserimento e rimozione.
    \item \textbf{Gestione della memoria}: Dopo ogni esperimento completo per un valore di $n$, la memoria viene liberata mediante garbage collection per evitare interferenze tra esperimenti successivi.
    \item \textbf{Raccolta dati robusta}: L'uso della mediana invece della media riduce l'impatto di misurazioni anomale e fornisce stime più affidabili.
\end{itemize}
\end{document}  